\chapter*{РЕАЛИЗАЦИЯ}
\addcontentsline{toc}{chapter}{РЕАЛИЗАЦИЯ}

\section*{Задание 1. Шифрование текста}

На листинге 1 приведена реализация функции шифрования текста шифром Виженера.

\includelistingpretty
    {vigenere_encrypt.cpp}
    {C++}
    {Шифрование текста шифром Виженера}

Функция encrypt выполняет шифрование текста с использованием ключа. Для каждого символа открытого текста вычисляется сдвиг по алфавиту, соответствующий текущему символу ключа. Ключ применяется циклически. Символы, не являющиеся буквами, остаются без изменений.

Программа была протестирована на тексте длиной 100000 символов. Результат работы программы с коротким ключом <<шифр>> (4 символа) представлен на рисунке~\ref{img:img1.png}.

\includeimage{img1.png}{f}{H}{0.9\textwidth}{Демонстрация шифрования с коротким ключом}

Результат работы программы с длинным ключом <<тавтологичность>> (15 символов) представлен на рисунке~\ref{img:img2.png}.

\includeimage{img2.png}{f}{H}{0.9\textwidth}{Демонстрация шифрования с длинным ключом}

\section*{Задание 2. Определение длины ключа}

На листинге 2 приведена реализация алгоритма Казиски для поиска повторяющихся n-грамм.

\includelistingpretty
    {kasiski.cpp}
    {C++}
    {Алгоритм Казиски}

Функция find\_key\_lengths выполняет поиск повторяющихся последовательностей символов в шифротексте и вычисляет расстояния между их позициями. НОД найденных расстояний используется для определения вероятных длин ключа.

На листинге 3 приведена реализация вычисления индекса совпадений.

\includelistingpretty
    {index_coincidence.cpp}
    {C++}
    {Вычисление индекса совпадений}

Функция calculate\_ic вычисляет индекс совпадений для текста. Функция find\_best\_key\_length перебирает возможные длины ключа и выбирает ту, при которой средний индекс совпадений подпоследовательностей наиболее близок к естественному для русского языка (0.0553).

Результат работы метода Казиски для криптограммы представлен на рисунке~\ref{img:img3.png}.

\includeimage{img3.png}{f}{H}{0.9\textwidth}{Демонстрация работы метода Казиски}

\section*{Задание 3. Взлом ключа}

На листинге 4 приведена реализация метода хи-квадрат для определения сдвига.

\includelistingpretty
    {chi_square.cpp}
    {C++}
    {Частотный анализ методом хи-квадрат}

Функция chi\_square вычисляет статистику $\chi^2$ для оценки совпадения частотного распределения текста с эталонным распределением русского языка при заданном сдвиге. Функция recover\_key последовательно определяет каждый символ ключа, выбирая сдвиг с минимальным значением $\chi^2$.

Результат работы программы по восстановлению ключа и расшифровке текста представлен на рисунке~\ref{img:img4.png}.

\includeimage{img4.png}{f}{H}{0.9\textwidth}{Демонстрация восстановления ключа и расшифровки}
